\documentclass[12pt, a4paper]{exam}
\usepackage[a4paper,top=2.3cm,bottom=2.8cm,left=2cm,right=2cm,headheight=15pt,headsep=0.6cm,footskip=1.2cm]{geometry}
\usepackage[T1]{fontenc}
\usepackage[utf8]{inputenc}
\usepackage{lmodern}
\usepackage{textcomp}
\usepackage[french]{babel}
\usepackage{xcolor}
\usepackage{listings}
\usepackage{graphicx}
\usepackage{array}
\usepackage{enumitem}
\usepackage{amsmath}
\usepackage{tikz}
\usetikzlibrary{positioning,shapes.multipart,arrows.meta}
\usepackage[most]{tcolorbox}

\definecolor{codegreen}{rgb}{0,0.6,0}
\definecolor{codegray}{rgb}{0.5,0.5,0.5}
\definecolor{codepurple}{rgb}{0.58,0,0.82}
\definecolor{backcolour}{rgb}{0.95,0.95,0.92}

\lstdefinestyle{mystyle}{
    backgroundcolor=\color{backcolour},
    commentstyle=\color{codegreen},
    keywordstyle=\color{magenta},
    numberstyle=\tiny\color{codegray},
    stringstyle=\color{codepurple},
    basicstyle=\ttfamily\footnotesize,
    breakatwhitespace=false,
    breaklines=true,
    postbreak=\mbox{\textcolor{codegray}{$\hookrightarrow$}\space},
    captionpos=b,
    keepspaces=true,
    columns=fullflexible,
    numbers=left,
    numbersep=5pt,
    showspaces=false,
    showstringspaces=false,
    showtabs=false,
    upquote=true,
    tabsize=4,
    xleftmargin=1.5em,
    framexleftmargin=0.3em,
    aboveskip=0.6em,
    belowskip=0.6em,
    extendedchars=true,
    literate={é}{{\'e}}1 {è}{{\`e}}1 {ê}{{\^e}}1 {ë}{{\"e}}1 {à}{{\`a}}1 {â}{{\^a}}1
             {ç}{{\c c}}1 {ù}{{\`u}}1 {û}{{\^u}}1 {ô}{{\^o}}1 {î}{{\^i}}1 {ï}{{\"i}}1
             {É}{{\'E}}1 {È}{{\`E}}1 {À}{{\`A}}1 {Ç}{{\c C}}1 {œ}{{\oe}}1
             {«}{{\guillemotleft}}1 {»}{{\guillemotright}}1 {→}{{$\rightarrow$}}1
}
\lstset{style=mystyle, language=Python}
\lstdefinestyle{sqlstyle}{style=mystyle, language=SQL, morekeywords={LIMIT,REFERENCES,IN,TEXT,FLOAT,DATE,VARCHAR,CHAR,REPLACE}, numbers=none, xleftmargin=0.3em, frame=lines}

% Code en ligne
\newcommand{\code}[1]{\texttt{#1}}

% Titre d'une question : numéro et intitulé, barème à droite
\newcommand{\Qt}[2]{\subsection*{#1\hfill{\normalfont\small #2}}}
% Rappel de la consigne en une ligne
\newcommand{\rappel}[1]{\noindent\textit{Ce qui est demandé :} #1\par\smallskip}

% La réponse attendue, mise en évidence
\newtcolorbox{reponse}[1][Réponse]{enhanced, breakable, colback=black!5, colframe=black, boxrule=0.9pt, arc=1mm,
  title={#1}, fonttitle=\bfseries\small, coltitle=black, colbacktitle=black!15, left=2mm, right=2mm, top=1mm, bottom=1mm,
  before skip=0.5em, after skip=0.7em}
% Erreur fréquente
\newtcolorbox{erreur}{enhanced, breakable, colback=white, colframe=black, boxrule=0.5pt, leftrule=3pt, arc=0mm,
  title={Erreur fréquente}, fonttitle=\bfseries\small, coltitle=black, colbacktitle=white, left=2mm, right=2mm, top=1mm, bottom=1mm,
  before skip=0.5em, after skip=0.7em}
% Encadré « À retenir »
\newtcolorbox{retenir}{enhanced, breakable, colback=white, colframe=black, boxrule=0.5pt, arc=1mm, boxsep=1mm,
  borderline={0.4pt}{1.2pt}{black, dashed},
  title={À retenir}, fonttitle=\bfseries\small, coltitle=black, colbacktitle=white, left=2mm, right=2mm, top=1mm, bottom=1mm,
  before skip=0.5em, after skip=0.7em}

\newcommand{\methode}{\noindent\textbf{Méthode.}\ }
\newcommand{\explication}{\noindent\textbf{Explication.}\ }

\setlist{itemsep=0.15em, topsep=0.3em}
\renewcommand{\arraystretch}{1.25}
\setlength{\parindent}{0pt}
\setlength{\parskip}{0.35em}

% Piles dessinées verticalement (sommet en haut) ; \suppr barre un élément supprimé
\newcommand{\suppr}[1]{\ooalign{#1\cr\hidewidth\rule[0.55ex]{1em}{1pt}\hidewidth\cr}}
\newcommand{\pilev}[2]{\begin{tikzpicture}[baseline=(titre.base)]
  \node[font=\scriptsize, align=center, text width=22mm] (titre) at (0,0) {#1};
  \foreach \v [count=\k] in {#2} {\node[draw, minimum width=8mm, minimum height=5mm, inner sep=0pt, font=\small] at (0,-0.25-\k*0.5) {\v};}
\end{tikzpicture}}
\setlength{\emergencystretch}{3em}
% Remarque sur l'énoncé (erreur ou ambiguïté du sujet)
\newtcolorbox{remarque}{enhanced, breakable, colback=white, colframe=black, boxrule=0.5pt, arc=0mm,
  borderline={0.4pt}{1.2pt}{black, dotted},
  title={Remarque sur l'énoncé}, fonttitle=\bfseries\small, coltitle=black, colbacktitle=black!8, left=2mm, right=2mm, top=1mm, bottom=1mm,
  before skip=0.5em, after skip=0.7em}

\title{Corrigé — CC1 Term NSI - Pile, File et rappels de 1er}
\author{Babash - NSI Tani Malandi et Johanne Bakalara}
\date{Sujet type bac de septembre 2025 (1 h 30, sur 20 points)}

\begin{document}
\pagestyle{headandfoot}
\firstpageheader{}{}{}
\runningheader{\small Corrigé}{}{\small Pile, file et rappels de 1re -- septembre 2025}
\runningheadrule
\firstpagefooter{}{\footnotesize Babash - NSI Tani Malandi et Johanne Bakalara \textperiodcentered{} CC BY-SA 4.0 \textperiodcentered{} page \thepage}{}
\runningfooter{}{\footnotesize Babash - NSI Tani Malandi et Johanne Bakalara \textperiodcentered{} CC BY-SA 4.0 \textperiodcentered{} page \thepage}{}

\maketitle
\thispagestyle{headandfoot}
\begin{center}\footnotesize\itshape D'après l'exercice 1 du sujet 22-NSIJ1G11 (baccalauréat 2022) pour les jours et les mois, et d'après le sujet 22-NSIJ2ME1 (baccalauréat 2022, Métropole, jour 2) pour la poussette.\end{center}

\begin{tcolorbox}[colback=white, colframe=black, boxrule=0.5pt, arc=1mm, left=2mm, right=2mm]
\small\textbf{Comment lire ce corrigé.} Pour chaque question : ce qui est demandé, puis la réponse attendue (encadrée), puis l'explication pas à pas. Essaie d'abord seul·e, puis compare ta démarche, pas seulement ton résultat. Tous les codes de ce corrigé ont été exécutés en Python 3.
Le barème est celui du sujet (11 questions, 20 points).
\end{tcolorbox}

\section*{Exercice 1 -- Dictionnaires et listes (10 points)}

On rappelle les données de l'exercice :
\begin{lstlisting}
jours = ["dimanche", "lundi", "mardi", "mercredi",
         "jeudi", "vendredi", "samedi"]
mois = {1 :("janvier",31), 2 :("février",28), 3 :("mars",31) ,
        4 :("avril",30), 5 :("mai",31), 6 :("juin",30) ,
        7 :("juillet",31), 8 :("aout",31), 9 :("septembre",30) ,
        10 :("octobre",31), 11 :("novembre",30), 12 :("décembre",31)}
\end{lstlisting}

\begin{retenir}
Dans une \textbf{liste}, chaque élément a un \textbf{indice} (sa position), et le premier indice est \textbf{0}.
\begin{center}\small
\begin{tabular}{|l|*{7}{c|}}\hline
indice & 0 & 1 & 2 & 3 & 4 & 5 & 6 \\\hline
élément & dimanche & lundi & mardi & mercredi & jeudi & vendredi & samedi \\\hline
indice négatif & $-7$ & $-6$ & $-5$ & $-4$ & $-3$ & $-2$ & $-1$ \\\hline
\end{tabular}
\end{center}
Dans un \textbf{dictionnaire}, on accède à une valeur par sa \textbf{clé} : ici, la clé est le numéro du mois et la valeur est un \textbf{tuple} (nom, nombre de jours).
\end{retenir}

\Qt{Question 1 -- Manipulation de liste}{2 points}

\textbf{a)} \rappel{obtenir l'élément \code{"lundi"} à partir de la liste \code{jours}.}
\begin{reponse}
\code{jours[1]} (cette expression vaut \code{"lundi"}).
\end{reponse}
\explication \code{"lundi"} est le deuxième élément de la liste. Comme on compte à partir de 0, le deuxième élément a l'indice 1. On écrit le nom de la liste, puis l'indice entre crochets.

\begin{erreur}
Écrire \code{jours[2]} en comptant à partir de 1 : on obtient \code{"mardi"}. En Python, le premier élément est toujours \code{jours[0]}.
\end{erreur}

\textbf{b)} \rappel{que renvoie \code{jours[-1]} ?}
\begin{reponse}
\code{jours[-1]} renvoie \code{"samedi"}, le dernier élément de la liste.
\end{reponse}
\explication Un indice négatif compte à partir de la fin : \code{-1} désigne le dernier élément, \code{-2} l'avant-dernier, etc. (voir la dernière ligne du tableau ci-dessus). C'est pratique quand on ne connaît pas la longueur de la liste : \code{jours[-1]} est la même chose que \mbox{\code{jours[len(jours) - 1]}}, c'est-à-dire \code{jours[6]}.

\Qt{Question 2 -- La fonction \code{jour\_suivant}}{2 points}
\rappel{écrire \code{jour\_suivant(jours, j)} qui renvoie le nom du jour qui suit le jour \code{j}.}

\methode
\begin{enumerate}
\item Trouver l'indice \code{i} du jour \code{j} avec \code{jours.index(j)}.
\item Le jour suivant est en général à l'indice \code{i + 1}.
\item Cas particulier : après \code{"samedi"} (indice 6, le dernier), on revient au début de la liste, à \code{"dimanche"} (indice 0).
\end{enumerate}

\begin{reponse}
\begin{lstlisting}
def jour_suivant(jours, j):
    i = jours.index(j)          # position du jour actuel
    if i == len(jours) - 1:     # j est le dernier jour (samedi)
        return jours[0]         # on revient à dimanche
    else:
        return jours[i + 1]     # sinon, c'est le jour d'après
\end{lstlisting}
\end{reponse}

\explication
\begin{itemize}
\item Ligne 2 : \code{jours.index("mardi")} vaut 2, donc \code{i} vaut 2.
\item Ligne 3 : \code{len(jours) - 1} vaut 6, l'indice du dernier élément. On teste si \code{j} est le dernier jour.
\item Lignes 4 et 6 : on \textbf{renvoie} (\code{return}) le nom du jour ; on ne l'affiche pas.
\end{itemize}

\begin{center}\small
\begin{tabular}{|l|c|c|l|}\hline
Appel & \code{i} & \code{i == 6} ? & Valeur renvoyée \\\hline
\code{jour\_suivant(jours, "mardi")} & 2 & non & \code{jours[3]}, soit \code{"mercredi"} \\\hline
\code{jour\_suivant(jours, "dimanche")} & 0 & non & \code{jours[1]}, soit \code{"lundi"} \\\hline
\code{jour\_suivant(jours, "samedi")} & 6 & oui & \code{jours[0]}, soit \code{"dimanche"} \\\hline
\end{tabular}
\end{center}

\textbf{Variante} avec le reste de la division euclidienne, noté \code{\%} en Python (on dit « modulo ») : \code{(6 + 1) \% 7} vaut 0, ce qui fait revenir au début de la liste.
\begin{lstlisting}
def jour_suivant(jours, j):
    return jours[(jours.index(j) + 1) % 7]
\end{lstlisting}

\begin{erreur}
Oublier le cas du samedi : \code{return jours[i + 1]} avec \code{i = 6} demande \code{jours[7]}, qui n'existe pas. Python affiche alors \code{IndexError: list index out of range} (« indice de liste hors limites »).
\end{erreur}

\Qt{Question 3 -- Manipulation de dictionnaire}{2 points}

\textbf{a)} \rappel{obtenir le nombre de jours du mois de mars à partir de \code{mois}.}
\begin{reponse}
\code{mois[3][1]} (cette expression vaut \code{31}).
\end{reponse}
\explication On lit l'expression de gauche à droite :
\begin{enumerate}
\item \code{mois[3]} : la valeur associée à la clé 3 (mars), c'est le tuple \code{("mars", 31)} ;
\item \code{("mars", 31)[1]} : l'élément d'indice 1 de ce tuple, c'est-à-dire \code{31}.
\end{enumerate}

\begin{erreur}
Écrire \code{mois["mars"]} : \code{"mars"} n'est pas une clé du dictionnaire (les clés sont les nombres de 1 à 12). Python affiche \code{KeyError: 'mars'} (« clé introuvable »).
\end{erreur}

\textbf{b)} \rappel{écrire \code{dans\_x\_mois(mois, numero\_mois, x)} qui renvoie le nom du mois qu'il sera \code{x} mois plus tard.}

\methode Les mois tournent en boucle : après 12 vient 1. Le reste de la division par 12 (\code{\% 12}) fait ce « tour du cadran », mais il donne un nombre entre 0 et 11. On procède donc ainsi :
\begin{enumerate}
\item on enlève 1 au numéro du mois pour se ramener à un nombre entre 0 et 11 ;
\item on ajoute \code{x}, puis on prend le reste de la division par 12 ;
\item on rajoute 1 pour revenir à un numéro entre 1 et 12 ;
\item on renvoie le nom du mois, à l'indice 0 du tuple.
\end{enumerate}

\begin{reponse}
\begin{lstlisting}
def dans_x_mois(mois, numero_mois, x):
    numero = (numero_mois - 1 + x) % 12 + 1   # numéro du mois futur
    return mois[numero][0]                    # son nom
\end{lstlisting}
\end{reponse}

\begin{center}\footnotesize
\begin{tabular}{|l|c|c|c|l|}\hline
Appel & \code{numero\_mois - 1 + x} & \code{\% 12} & \code{+ 1} & Valeur renvoyée \\\hline
\code{dans\_x\_mois(mois, 4, 5)} & $3 + 5 = 8$ & 8 & 9 & \code{mois[9][0]}, soit \code{"septembre"} \\\hline
\code{dans\_x\_mois(mois, 10, 3)} & $9 + 3 = 12$ & 0 & 1 & \code{mois[1][0]}, soit \code{"janvier"} \\\hline
\end{tabular}
\end{center}

\textbf{Variante} sans modulo : on avance d'un mois, \code{x} fois de suite.
\begin{lstlisting}
def dans_x_mois(mois, numero_mois, x):
    numero = numero_mois
    for k in range(x):        # on répète x fois
        if numero == 12:
            numero = 1        # après décembre : janvier
        else:
            numero = numero + 1
    return mois[numero][0]
\end{lstlisting}
Trace pour \code{dans\_x\_mois(mois, 10, 3)} : \code{numero} vaut 10 au départ, puis 11 (tour \code{k = 0}), 12 (tour \code{k = 1}), 1 (tour \code{k = 2}). La fonction renvoie \code{mois[1][0]}, soit \code{"janvier"}.

\begin{erreur}
Écrire \code{(numero\_mois + x) \% 12} sans le « $-1$ » et le « $+1$ » : pour décembre, on obtient 0. Par exemple \code{dans\_x\_mois(mois, 10, 2)} chercherait \code{mois[0]}, qui n'existe pas : \code{KeyError: 0}, au lieu de \code{"décembre"}.
\end{erreur}

\Qt{Question 4 -- Les dates sous forme de tuple}{3 points}

\textbf{a)} \rappel{on pose \code{date = ("samedi", 21, 10, 1995)}. Que renvoie l'expression \code{mois[date[2]][1]} ?}
\begin{reponse}
\code{mois[date[2]][1]} renvoie \code{31}.
\end{reponse}
\explication On calcule de l'intérieur vers l'extérieur :
\begin{enumerate}
\item \code{date[2]} : l'élément d'indice 2 du tuple \code{date}, c'est \code{10} (le numéro du mois) ;
\item \code{mois[10]} : le tuple \code{("octobre", 31)} ;
\item \code{("octobre", 31)[1]} : \code{31}, le nombre de jours du mois d'octobre.
\end{enumerate}

\textbf{b)} \rappel{que se passe-t-il si on exécute \code{date[0] = "dimanche"} ? Pourquoi ?}
\begin{reponse}
L'instruction provoque une erreur : \code{TypeError: 'tuple' object does not support item assignment}. Un tuple n'est pas modifiable.
\end{reponse}
\explication Le message signifie « un objet de type tuple n'accepte pas l'affectation d'un élément ». Une liste est \textbf{mutable} : on peut changer ses éléments après sa création. Un tuple est \textbf{non mutable} (on dit aussi \textbf{immuable}) : une fois créé, on ne peut ni changer, ni ajouter, ni retirer un élément. Pour « changer » une date, on fabrique un \textbf{nouveau} tuple, par exemple \code{date = ("dimanche", 22, 10, 1995)}.

\textbf{c)} \rappel{compléter \code{changement\_mois(numero\_mois, annee)}, qui renvoie le tuple (mois suivant, année).}
\begin{reponse}
\begin{lstlisting}
def changement_mois(numero_mois, annee):
    if numero_mois == 12:
        return (1, annee + 1)
    else:
        return (numero_mois + 1, annee)
\end{lstlisting}
\end{reponse}
\explication
\begin{itemize}
\item Si on est en décembre (mois 12), le mois suivant est janvier (1) \textbf{et} on change d'année : \code{annee + 1}.
\item Sinon, on passe au mois suivant (\code{numero\_mois + 1}) et l'année ne change pas.
\end{itemize}
Vérification : \code{changement\_mois(1, 2025)} renvoie \code{(2, 2025)} et \code{changement\_mois(12, 1900)} renvoie \code{(1, 1901)}.

\Qt{Question 5 -- La fonction \code{date\_suivante}}{1 point}
\rappel{écrire \code{date\_suivante(date)} qui renvoie le tuple de la date du lendemain.}

\methode
\begin{enumerate}
\item Récupérer les quatre éléments du tuple (nom du jour, numéro du jour, mois, année).
\item Le nom du jour change toujours : on utilise \code{jour\_suivant}.
\item Regarder si on est le dernier jour du mois : on compare le numéro du jour avec \code{mois[numero\_mois][1]}.
\item Si oui, le lendemain est le 1\textsuperscript{er} du mois suivant : on utilise \code{changement\_mois}. Sinon, on ajoute 1 au numéro du jour.
\end{enumerate}

\begin{reponse}
\begin{lstlisting}
def date_suivante(date):
    nom_jour, numero_jour, numero_mois, annee = date
    nouveau_nom = jour_suivant(jours, nom_jour)
    if numero_jour == mois[numero_mois][1]:      # dernier jour du mois
        nouveau_mois, nouvelle_annee = changement_mois(numero_mois, annee)
        return (nouveau_nom, 1, nouveau_mois, nouvelle_annee)
    else:
        return (nouveau_nom, numero_jour + 1, numero_mois, annee)
\end{lstlisting}
\end{reponse}

\explication
\begin{itemize}
\item Ligne 2 : on « déballe » le tuple en quatre variables, dans l'ordre (\code{nom\_jour}, \code{numero\_jour}, \code{numero\_mois}, \code{annee}). On peut aussi écrire \code{nom\_jour = date[0]}, \code{numero\_jour = date[1]}, etc.
\item Ligne 5 : \code{changement\_mois} renvoie un tuple de deux valeurs, qu'on range dans deux variables.
\item Lignes 6 et 8 : on renvoie un \textbf{nouveau} tuple (on ne peut pas modifier \code{date}, voir la question 4~b).
\end{itemize}

\begin{center}\small
\begin{tabular}{|l|c|l|}\hline
Date de départ & Dernier jour du mois ? & Date renvoyée \\\hline
\code{("samedi", 21, 10, 1995)} & non ($21 \neq 31$) & \code{("dimanche", 22, 10, 1995)} \\\hline
\code{("mardi", 31, 10, 1995)} & oui ($31 = 31$) & \code{("mercredi", 1, 11, 1995)} \\\hline
\code{("dimanche", 31, 12, 1995)} & oui & \code{("lundi", 1, 1, 1996)} \\\hline
\code{("mardi", 28, 2, 2023)} & oui (février : 28) & \code{("mercredi", 1, 3, 2023)} \\\hline
\end{tabular}
\end{center}

\begin{erreur}
Tester \code{numero\_jour == 31} pour tous les mois : le 30 avril serait suivi d'un « 31 avril ». Il faut lire le nombre de jours du mois dans le dictionnaire \code{mois}.
\end{erreur}

\section*{Exercice 2 -- Piles : la poussette (10 points)}

\begin{retenir}
Une \textbf{pile} est une structure de données où l'on ajoute (\textbf{empiler}) et retire (\textbf{dépiler}) les éléments \textbf{au même bout}, le \textbf{sommet}. Le dernier élément ajouté est le premier retiré.
Une \textbf{file} ajoute à un bout et retire à l'autre bout : le premier élément ajouté est le premier retiré.
Dans ce corrigé, une pile s'écrit du haut vers le bas : « 4 9 8 7 4 2 » veut dire que 4 est au sommet et 2 tout en bas.
\end{retenir}

\Qt{Question 6 -- Question de cours}{2 points}

\textbf{a)} \rappel{que veut dire LIFO ? Pile ou file ?}
\begin{reponse}
LIFO veut dire \textit{Last In, First Out}, c'est-à-dire « dernier entré, premier sorti ». Cela caractérise une \textbf{pile}.
\end{reponse}
\explication Dans une pile, on retire toujours l'élément du sommet, qui est le dernier ajouté. Une file, au contraire, est FIFO : \textit{First In, First Out}, « premier entré, premier sorti ».

\textbf{b)} \rappel{un exemple de file et un exemple de pile dans la vie courante.}
\begin{reponse}
\textbf{Pile :} une pile d'assiettes (on pose et on prend l'assiette du dessus) ; le bouton « Annuler » d'un logiciel, qui défait d'abord la dernière action.\\
\textbf{File :} la file d'attente des voitures à la barge (la première arrivée embarque la première) ; la file d'attente des documents envoyés à une imprimante.
\end{reponse}
\explication Pour reconnaître une pile ou une file, pose-toi la question : « qui sort en premier ? ». Si c'est le dernier arrivé, c'est une pile ; si c'est le premier arrivé, c'est une file.

\Qt{Question 7 -- Réduire une pile à la main}{2 points}

\begin{retenir}
\textbf{Comment on parcourt la pile} (c'est ce que montre l'exemple de l'énoncé, et ce que fait le code des questions 8 et 9) :
\begin{enumerate}
\item on regarde les trois éléments du haut ;
\item si celui du haut et celui du bas du triplet ont la même parité, on supprime celui du milieu ;
\item on met de côté l'élément du haut, puis on recommence avec les trois éléments suivants ;
\item quand il reste moins de trois éléments, le parcours est fini : on remet les éléments mis de côté ;
\item on refait des parcours tant qu'un parcours supprime au moins un élément.
\end{enumerate}
\end{retenir}

\textbf{a)} \rappel{donner les étapes de réduction de la pile 4 9 8 7 4 2 (du haut vers le bas).}
\begin{reponse}
\begin{tabular}{@{}lll@{}}
1\textsuperscript{er} parcours : & 4 9 8 7 4 2 $\rightarrow$ 4 8 7 4 2 $\rightarrow$ 4 8 4 2 & (on supprime 9, puis 7)\\
2\textsuperscript{e} parcours : & 4 8 4 2 $\rightarrow$ 4 4 2 & (on supprime 8)\\
3\textsuperscript{e} parcours : & 4 4 2 $\rightarrow$ 4 2 & (on supprime le 4 du milieu)\\
\end{tabular}\\
La pile finale est \textbf{4 2} : elle a deux éléments, la partie est gagnante.
\end{reponse}

\explication On applique la méthode. « P » veut dire pair, « I » impair.
\begin{center}\small
\begin{tabular}{|c|l|c|c|l|}\hline
Parcours & Triplet examiné & Haut / bas & Action & Pile après \\\hline
1 & 4 9 8 & P / P & on supprime 9 & 4 8 7 4 2 \\
1 & 8 7 4 (on a mis 4 de côté) & P / P & on supprime 7 & 4 8 4 2 \\
1 & il reste 4 2 sous le 8 : fin & & & 4 8 4 2 \\\hline
2 & 4 8 4 & P / P & on supprime 8 & 4 4 2 \\
2 & il reste 4 2 : fin & & & 4 4 2 \\\hline
3 & 4 4 2 & P / P & on supprime 4 (milieu) & 4 2 \\\hline
4 & moins de 3 éléments & & plus rien à réduire & 4 2 \\\hline
\end{tabular}
\end{center}

\begin{center}
\pilev{départ}{4,\suppr{9},8,7,4,2}\quad
\pilev{après 9}{4,8,\suppr{7},4,2}\quad
\pilev{fin du\\1\textsuperscript{er} parcours}{4,\suppr{8},4,2}\quad
\pilev{fin du\\2\textsuperscript{e} parcours}{4,\suppr{4},2}\quad
\pilev{pile finale}{4,2}\\[0.3em]
{\small Sur chaque dessin, l'élément barré est celui qui va être supprimé à l'étape suivante.}
\end{center}

\textbf{b)} \rappel{parmi les piles A, B et C, donner celle qui est gagnante.}

On lit chaque pile du haut vers le bas sur le dessin : A = 5 4 5 4 2 1 ; B = 4 5 4 9 2 0 ; C = 3 4 8 7 6 1.
\begin{reponse}
La pile gagnante est la \textbf{pile B} : elle se réduit à \textbf{4 0}, soit deux éléments.
\end{reponse}
\explication
\begin{itemize}
\item \textbf{Pile B} :
1\textsuperscript{er} parcours : 4 5 4 (P / P) on supprime 5 ; 4 9 2 (P / P) on supprime 9 ; il ne reste que 2 0 : on obtient 4 4 2 0.
2\textsuperscript{e} parcours : 4 4 2 (P / P) on supprime le 4 du milieu : 4 2 0.
3\textsuperscript{e} parcours : 4 2 0 (P / P) on supprime 2 : \textbf{4 0}. Deux éléments : gagnante.
\item \textbf{Pile A} : 5 4 5 (I / I) on supprime 4 : 5 5 4 2 1. Ensuite, aucun triplet ne se réduit (5 5 4 : I / P ; 5 4 2 : I / P ; 4 2 1 : P / I). La pile finale a 5 éléments : perdante.
\item \textbf{Pile C} : 3 4 8 (I / P) rien ; 4 8 7 (P / I) rien ; 8 7 6 (P / P) on supprime 7 : 3 4 8 6 1. Puis 4 8 6 (P / P) on supprime 8 : 3 4 6 1. Plus rien ne se réduit : 4 éléments, perdante.
\end{itemize}

\Qt{Question 8 -- La fonction \code{reduire\_triplet\_au\_sommet}}{1 point}
\rappel{compléter la fonction qui supprime l'élément central des trois éléments du haut quand le haut et le bas ont la même parité.}
\begin{reponse}
\begin{lstlisting}
def reduire_triplet_au_sommet(p):
    a = depiler(p)
    b = depiler(p)
    c = sommet(p)
    if a % 2 != c % 2:
        empiler(p, b)
    empiler(p, a)
\end{lstlisting}
\end{reponse}
\explication
\begin{itemize}
\item Lignes 2 à 4 : on retire \code{a} (le haut) et \code{b} (le milieu). On \textbf{lit} \code{c} avec \code{sommet}, sans le retirer : il reste dans la pile.
\item Ligne 5 : \code{a \% 2} vaut 0 si \code{a} est pair et 1 s'il est impair. La condition \code{a \% 2 != c \% 2} est donc vraie quand \code{a} et \code{c} n'ont \textbf{pas} la même parité : il n'y a pas de réduction, on remet \code{b}.
\item Ligne 7 : dans tous les cas, on remet \code{a} au sommet. Quand \code{a} et \code{c} ont la même parité, \code{b} n'a pas été remis : il est supprimé.
\item On remet \code{b} \textbf{avant} \code{a} : le dernier empilé se retrouve au sommet, donc \code{a} est bien au-dessus de \code{b}, comme au départ.
\end{itemize}
Exemples : avec 1 0 3, on a \code{a = 1}, \code{b = 0}, \code{c = 3} : même parité, on ne remet que 1, la pile devient 1 3. Avec 1 0 8, les parités diffèrent : on remet 0 puis 1, la pile reste 1 0 8.

\begin{erreur}
Empiler \code{a} puis \code{b} : l'ordre est inversé. Avec 1 0 8, on obtiendrait 0 1 8 au lieu de 1 0 8.
\end{erreur}

\Qt{Question 9 -- Parcourir la pile en la réduisant}{3 points}

\textbf{a)} \rappel{la taille minimale d'une pile réductible.}
\begin{reponse}
Il faut au moins \textbf{3 éléments}.
\end{reponse}
\explication La règle porte sur un triplet : il faut un élément en haut, un au milieu (celui qu'on peut supprimer) et un en bas. Avec deux éléments, il n'y a pas d'élément du milieu.

\textbf{b)} \rappel{compléter \code{parcourir\_pile\_en\_reduisant}.}
\begin{reponse}
\begin{lstlisting}
def parcourir_pile_en_reduisant(p):
    q = creer_pile_vide()
    while taille(p) >= 3:
        reduire_triplet_au_sommet(p)
        e = depiler(p)
        empiler(q, e)
    while not est_vide(q):
        e = depiler(q)
        empiler(p, e)
    return p
\end{lstlisting}
\end{reponse}
\explication
\begin{itemize}
\item Ligne 2 : \code{q} est une pile auxiliaire (une pile « de côté »). Elle reçoit les éléments déjà examinés.
\item Lignes 3 à 6 : tant qu'il reste un triplet, on le réduit si possible, puis on retire l'élément du sommet de \code{p} pour le poser sur \code{q}. Le triplet suivant commence donc un cran plus bas.
\item Lignes 7 à 9 : à la fin, \code{p} contient ses deux derniers éléments. On reprend un à un les éléments de \code{q} pour les remettre sur \code{p}. Comme \code{q} est une pile, le premier élément repris est le dernier posé : chaque élément retrouve sa place d'origine.
\end{itemize}

Trace pour la pile 4 9 8 7 4 2 (premier parcours de la question 7) :
\begin{center}\small
\begin{tabular}{|l|l|l|l|l|}\hline
Étape & \code{p} avant & Réduction & \code{e} & \code{p} et \code{q} après \\\hline
tour 1 & 4 9 8 7 4 2 & 4 / 8 pairs : 9 supprimé & 4 & \code{p} = 8 7 4 2 ; \code{q} = 4 \\
tour 2 & 8 7 4 2 & 8 / 4 pairs : 7 supprimé & 8 & \code{p} = 4 2 ; \code{q} = 8 4 \\\hline
\multicolumn{5}{|l|}{\code{taille(p)} vaut 2 : on sort de la première boucle.} \\\hline
retour 1 & 4 2 & & 8 & \code{p} = 8 4 2 ; \code{q} = 4 \\
retour 2 & 8 4 2 & & 4 & \code{p} = 4 8 4 2 ; \code{q} vide \\\hline
\end{tabular}
\end{center}
La fonction renvoie la pile 4 8 4 2, comme à la question 7~a.

\begin{erreur}
Écrire \code{while taille(p) > 3} : le dernier triplet n'est jamais examiné. Par exemple, la pile 1 0 3 resterait 1 0 3 alors qu'il faut supprimer le 0.
\end{erreur}

\Qt{Question 10 -- La fonction \code{est\_gagnant}}{1 point}
\rappel{renvoyer \code{True} si la pile finale a exactement deux éléments, \code{False} sinon.}
\begin{reponse}
\begin{lstlisting}
def est_gagnant(pile_finale):
    return taille(pile_finale) == 2
\end{lstlisting}
\end{reponse}
\explication \code{taille(pile\_finale) == 2} est une \textbf{comparaison} : elle vaut déjà \code{True} ou \code{False}. On peut donc la renvoyer directement. Variante plus longue, tout aussi juste :
\begin{lstlisting}
def est_gagnant(pile_finale):
    if taille(pile_finale) == 2:
        return True
    else:
        return False
\end{lstlisting}

\begin{erreur}
Confondre \code{=} et \code{==} : \code{=} range une valeur dans une variable (affectation), \code{==} compare deux valeurs. \code{return taille(pile\_finale) = 2} provoque une \code{SyntaxError}.
\end{erreur}

\Qt{Question 11 -- Taille renvoyée par \code{jouer}}{1 point}
\rappel{avec la pile 4 5 4 9 2 0, que renvoie \code{taille(jouer(p))} ?}
\begin{reponse}
\code{taille(jouer(p))} renvoie \textbf{2} : la pile finale est 4 0.
\end{reponse}
\explication La fonction \code{jouer} garde dans \code{q} une copie de la pile \textbf{avant} chaque parcours (\code{copier\_pile} copie la pile sans la modifier). Elle s'arrête quand un parcours ne change plus la taille, c'est-à-dire quand plus rien ne se réduit. C'est la pile B de la question 7~b.
\begin{center}\small
\begin{tabular}{|l|c|l|c|l|}\hline
Moment & \code{taille(q)} & \code{p} après le parcours & \code{taille(p)} & Tailles différentes ? \\\hline
avant la boucle & 6 & 4 4 2 0 & 4 & oui : on continue \\
tour 1 & 4 & 4 2 0 & 3 & oui : on continue \\
tour 2 & 3 & 4 0 & 2 & oui : on continue \\
tour 3 & 2 & 4 0 & 2 & non : on s'arrête \\\hline
\end{tabular}
\end{center}
La fonction renvoie la pile 4 0, de taille 2 : la partie est gagnante.

\end{document}
